% LTeX: language=de

\subsubsection{ln-Funktion}

\begin{note}
   \begin{itemize}
    \item allgemein gilt: \qquad {
        \begin{minipage}[t]{0.8\linewidth}
            \vspace{-3.0ex}
            \begin{align*}
                b &= a^{x} & \textcolor{myg}{9} &\textcolor{myg}{= 3^{x}} \\[1.5ex]
                x &= \log_{a}(b) & \textcolor{myg}{x} &\textcolor{myg}{= \log_{3}(9)}
            \end{align*}\\[-1.0ex]
            $\boxed{y = e^{x} \quad \Leftrightarrow \quad x = \log_{e}(y) = \ln(y)}$\\
        \end{minipage}
        }
    \item $\rightarrow ln$ bestimmt den Exponenten einer e-Funktion\\[1.0ex]
        $f(x) \underbrace{\coloneqq}_{\textcolor{myg}{\text{\footnotesize wird definiert}}} e^{x}$ \qquad \qquad $ x = ln(f(x)) = ln(e^{x}) = x$\\[1.0ex]
    \item $\rightarrow ln(x)$ ist die \textbf{Umkehrfunktion} von $e^{x}$
    \item Definitions- und Wertemenge: \\{
        \begin{minipage}[t]{0.60\linewidth}
            \raisebox{-\height}{
            \begin{tikzpicture}[
                    declare function={
                        w_1(\x) = (\x*\xScale)*(1/\yScale);
                        f_f(\x) = (exp(\x*\xScale))*(1/\yScale);
                        f_g(\x) = (ln(\x*\xScale))*(1/\yScale);
                    }
                ]

                % Set variables
                \pgfmathsetmacro{\axisPosWidth}{4};
                \pgfmathsetmacro{\axisNegWidth}{-3};
                \pgfmathsetmacro{\axisPosHeight}{5};
                \pgfmathsetmacro{\axisNegHeight}{-3};
                \pgfmathsetmacro{\xIncrement}{1.0};
                \pgfmathsetmacro{\yIncrement}{1.0};
                \pgfmathsetmacro{\xScale}{1};
                \pgfmathsetmacro{\yScale}{1};
                \pgfmathsetmacro{\gridxIncrement}{0.5};
                \pgfmathsetmacro{\gridyIncrement}{0.5};

                % Axis/grid limits
                \pgfmathsetmacro{\xUpperLimitAxis}{\xIncrement * floor(\axisPosWidth / \xIncrement)}
                \pgfmathsetmacro{\xLowerLimitAxis}{\xIncrement * ceil(\axisNegWidth / \xIncrement)}
                \pgfmathsetmacro{\yUpperLimitAxis}{\yIncrement * floor(\axisPosHeight /\yIncrement)}
                \pgfmathsetmacro{\yLowerLimitAxis}{\yIncrement * ceil(\axisNegHeight / \yIncrement)}

                \pgfmathsetmacro{\xUpperLimitGrid}{\gridxIncrement * floor(\axisPosWidth / \gridxIncrement)}
                \pgfmathsetmacro{\xLowerLimitGrid}{\gridxIncrement * ceil(\axisNegWidth / \gridxIncrement)}
                \pgfmathsetmacro{\yUpperLimitGrid}{\gridyIncrement * floor(\axisPosHeight / \gridyIncrement)}
                \pgfmathsetmacro{\yLowerLimitGrid}{\gridyIncrement * ceil(\axisNegHeight / \gridyIncrement)}

                % Axis environment (PGFPlots)
                \begin{axis}[
                    xmin=\xLowerLimitGrid, xmax=\xUpperLimitGrid,
                    ymin=\yLowerLimitGrid, ymax=\yUpperLimitGrid,
                    axis lines=middle,
                    axis line style={->, >=stealth},
                    grid=both,
                    xlabel={$x$},
                    ylabel={$y$},
                    xlabel style={at={(axis description cs:1,0.45)}, anchor=north},
                    ylabel style={at={(axis description cs:0.45,1)}, anchor=south},
                    xtick={\xLowerLimitAxis,\xLowerLimitAxis+1,...,\xUpperLimitAxis},
                    ytick={\yLowerLimitAxis,\yLowerLimitAxis+1,...,\yUpperLimitAxis},
                    xticklabel={\pgfmathparse{\tick*\xScale}\pgfmathprintnumber{\pgfmathresult}},
                    yticklabel={\pgfmathparse{\tick*\yScale}\pgfmathprintnumber{\pgfmathresult}},
                    tick label style={font=\scriptsize},
                    x=1cm,
                    y=1cm,
                    minor tick num=1,
                    scale only axis,
                    grid=both,
                    restrict y to domain=\yLowerLimitGrid:\yUpperLimitGrid,
                ]

                % G_f
                \addplot[thick, myb, samples=1000, domain=\xLowerLimitGrid:\xUpperLimitGrid]
                    {f_f(x)};
                \node[anchor=north] at (axis cs:{1*(1/\xScale)},{2*(1/\yScale)}) {\textcolor{myb}{$G_{e^{x}}$}};

                % G_g
                \addplot[thick, myr, samples=1000, domain=0.001:\xUpperLimitGrid]
                    {f_g(x)};
                \node[anchor=north] at (axis cs:{4.0*(1/\xScale)},{1.0*(1/\yScale)}) {\textcolor{myr}{$G_{ln(x)}$}};

                % P1
                \addplot[only marks, mark=x, mark size=3pt] coordinates {(0,1)};
                \node[anchor=south east] at (axis cs:0.0,1.0) {$P_{1}(0,1)$};

                % P2
                \addplot[only marks, mark=x, mark size=3pt] coordinates {(1,0)};
                \node[anchor=north west] at (axis cs:1,-0.25) {$P_{2}(1,0)$};

                % Winkelhalbierende
                \addplot[thick, myg, samples=1000, domain=\xLowerLimitGrid:0.96*\xUpperLimitGrid]
                    {w_1(x)};
                \end{axis}
            \end{tikzpicture}
            }\\[3.0ex]   
        \end{minipage}
        \begin{minipage}[t]{0.30\linewidth}
            \vspace{-2.0ex}
            \begin{align*}
                D_{e^{x}} &= \mathbb{R} &= W_{ln(x)}\\[1.5ex]
                W_{e^{x}} &= \mathbb{R}^{+} = \left] \ 0; \infty \right] &= D_{ln(x)}
            \end{align*}
        \end{minipage}
        }
    \item Eigenschaften von $ln(x)$: \\{
        \vspace{-2.0ex}
        \begin{itemize}
            \item $G_{ln(x)}$ ist sms \\[0.50ex]
            \item $\lim\limits_{x \to 0^+} ln(x) = -\infty$ \qquad $\lim\limits_{x \to \infty} ln(x) = \infty$\\[0.50ex]
            \item $\highlight[yellow]{ln(1) = 0}$, da $e^{0} = 1$
        \end{itemize}
        }
   \end{itemize} 
\end{note}

\begin{Example}{ln-Funktion}{}
    \begin{enumerate}
        \item [a)]{
            $f(x) = ln(x^2 -4);$ \qquad $w(x) = x^2 -4 \quad \textcolor{myg}{\text{\} \quad \footnotesize Argument der ln-Funktion} \ (> 0!)}$\\[-2.0ex]
            \begin{itemize}[leftmargin=*]
                \item Definitionsmenge $D_{f}$:\\{
                    \vspace{-2.0ex}
                    \begin{align*}
                        w(x) &= 0 &&\\[2.0ex]
                        x^2 - 4 &= 0 &&| \quad + 4\\[1.5ex]
                        x^2 &= 4 &&| \quad \sqrt[2]{\phantom{x}}\\[1.5ex]
                        x_{1} = -2 \text{\textit{(e)}}; &\quad x_{2} = 2 \text{\textit{(e)}} &&
                    \end{align*}\\
                    \begin{tabularx}{0.8\linewidth}{c|l|l|l|l|l}
                        & $-\infty < x < -2$ & $x = -2$ & $-2 < x < 2$ & $x = 2$ & $2 < x < \infty$ \\ \hline
                        $w(x)$ & $+$ & $0$ & $-$ & $0$ & $+$ \\ \hline
                        $f(x)$ & $+$ & n. d. & n. d. & n. d. & $+$ \\                   
                    \end{tabularx}\\[3.0ex]
                    $\Rightarrow \uuline{D_{f} = \left] \ -\infty; -2 \ \right[ \quad \cup \quad \left] \ 2; \infty \ \right[}$\\
                    }
                \item Nullstellen:\\[2.0ex]{
                    $f(x) = 0$ \qquad $\Rightarrow$ \quad $w(x) = 1$\\[-1.0ex]
                    \begin{align*}
                        w(x) &= 1 \\[1.5ex]
                        x^2 - 4 &= 1 &&| \quad + 4\\[1.5ex]
                        x^2 &= 5 &&| \quad \pm\sqrt[2]{\phantom{x}}\\[1.5ex]
                        x_{1} = -\sqrt[2]{5} \quad \text{\textit{(e)}}; &\quad x_{2} = \sqrt[2]{5} \quad \text{\textit{(e)}} &&\\[1.5ex]
                        x_{1} \in D_{f}; &\quad x_{2} \in D_{f} &&
                    \end{align*}
                    }
                \item Verhalten am Rand von $D_{f}$:\\{
                    \vspace{-2.0ex}
                    \begin{align*}
                        \lim\limits_{x \to -\infty} ln(\overbrace{x^2 - 4}^{\textcolor{myg}{\infty}}) &= +\infty; \qquad &\lim\limits_{x \to \infty} ln(\overbrace{x^2 - 4}^{\textcolor{myg}{\infty}}) &= +\infty &&\\[2.0ex]
                        \lim\limits_{x \to -2^-} ln(\overbrace{x^2 - 4}^{\textcolor{myg}{0^{+}}}) &= -\infty; \qquad &\lim\limits_{x \to 2^+} ln(\overbrace{x^2 - 4}^{\textcolor{myg}{0^{+}}}) &= -\infty &&
                    \end{align*}
                    }
            \end{itemize}
            }
    \end{enumerate}
\end{Example}